\documentclass[11pt,a4paper]{article}

\usepackage[margin=1in]{geometry}
\usepackage{amsmath,amssymb,mathtools}
\usepackage{enumitem}

\newcommand{\StartE}{\mathit{StartE}}
\newcommand{\EndE}{\mathit{EndE}}
\newcommand{\GType}{\mathit{GType}}
\newcommand{\gtype}{\mathit{gtype}}
\newcommand{\ActB}{\mathit{Act}_{B}}
\newcommand{\node}{\mathit{node}}
\newcommand{\InF}{\mathit{InF}}
\newcommand{\OutF}{\mathit{OutF}}
\newcommand{\ArrB}{\mathit{Arr}_{B}}
\newcommand{\JoinPre}{\mathit{JoinPre}_{B}}
\newcommand{\SelectedB}{\mathit{Selected}_{B}}
\newcommand{\ConsumeB}{\mathit{Consume}_{B}}
\newcommand{\ProduceB}{\mathit{Produce}_{B}}

\title{Core BPMN Process-State Semantics}
\author{}
\date{}

\begin{document}
\maketitle

\section{BPMN Process Structure}

A BPMN process is a tuple
\[
B = (T, F)
\]
where
\[
T = E \uplus A \uplus G
\]
is a finite set of process transitions, consisting of Events $E$, Activities $A$, and Gateways $G$. Events are partitioned into Start Events and End Events:
\[
E = \StartE \uplus \EndE.
\]
The sequence-flow relation is
\[
F \subseteq T \times T.
\]

\section{Predecessors and Successors}

For each transition $t \in T$, its predecessor and successor sets are defined by
\[
{}^{\bullet}t = \{ t' \in T \mid (t',t) \in F \}
\]
and
\[
t^{\bullet} = \{ t' \in T \mid (t,t') \in F \}.
\]

\section{Structural Well-Formedness}

The BPMN profile satisfies the following structural constraints.

\paragraph{Start Events.}
\[
\forall s \in \StartE:\quad
{}^{\bullet}s = \varnothing
\;\land\;
|s^{\bullet}| = 1
\tag{WF-S}
\]

\paragraph{End Events.}
\[
\forall e \in \EndE:\quad
|{}^{\bullet}e| \geq 1
\;\land\;
e^{\bullet} = \varnothing
\tag{WF-E}
\]

\paragraph{Activities.}
\[
\forall a \in A:\quad
|{}^{\bullet}a| \geq 1
\;\land\;
|a^{\bullet}| = 1
\tag{WF-A}
\]

\paragraph{Gateways.}
\[
\forall g \in G:\quad
|{}^{\bullet}g| \geq 1
\;\land\;
|g^{\bullet}| \geq 1
\tag{WF-G1}
\]
A gateway with exactly one incoming and one outgoing sequence flow is excluded:
\[
\forall g \in G:\quad
|{}^{\bullet}g| > 1
\;\lor\;
|g^{\bullet}| > 1
\tag{WF-G2}
\]

\section{Gateway Types}

The gateway types considered in the current profile are
\[
\GType = \{ \mathit{Exclusive}, \mathit{Inclusive}, \mathit{Parallel} \}.
\]
Each gateway has exactly one type:
\[
\gtype : G \rightarrow \GType.
\]
The operational effect of each gateway type is captured later by its join precondition and outgoing-flow selection rule.




\section{Extended Flow Domain}

Since Start Events have no incoming sequence flow, let $\bot \notin T$ be a distinguished process-entry symbol and define
\[
F^{0} = \{ (\bot, s) \mid s \in \StartE \}.
\]
The flow domain used by the process-state semantics is
\[
F_B = F \cup F^{0}.
\]

A BPMN process state is denoted by
\[
m(\sigma_F) \mid \sigma_F \subseteq F_B.
\]

The initial configuration is
\[
\mathit{Initial}_B(m(\sigma_F))
\iff
\sigma_F = F^{0}.
\]

A process state is final iff all active flows lead to End Events:
\[
\mathit{Final}_B(m(\sigma_F))
\iff
\forall f \in \sigma_F:\ \mathit{target}(f) \in \EndE.
\]

\section{State-Transition Semantics}

For each flow $f = (u, t) \in F_B$, define
\[
\mathit{source}(f) = u,
\qquad
\mathit{target}(f) = t.
\]

Every flow in the extended flow domain is executable, including a flow whose
target is an End Event:
\[
F_B^{exec}
=F_B.
\]

The execution of $t$ induces a transition between two flow configurations:
\[
\sigma_F \xrightarrow{\mathit{target}(f)}_{B} \sigma_F',\quad f \in \sigma_F.
\]

The resulting configuration $\sigma_F'$ is determined by the successor function:
\[
\frac{
f \in \sigma_F \cap F_B^{exec}
\qquad
\mathit{Pass}(f) \subseteq \sigma_F
}{
\sigma_F \xrightarrow{\mathit{target}(f)} \sigma_F' :
\sigma_F' = \bigl(\sigma_F \setminus \mathit{Pass}(f)\bigr)
\cup
\mathit{Next}(f)
}.
\tag{EXEC}
\]

\section{Flow Passing and Successor Selection}

Besides the predecessor and successor node sets, define the incident flow sets
\[
\InF(t) = \{ (u, t) \mid u \in T \land (u, t) \in F \}
\qquad\text{and}\qquad
\OutF(t) = \{ (t, v) \mid v \in T \land (t, v) \in F \}.
\]

Let $f \in F_B^{exec}$ and $t = \mathit{target}(f)$.

For every nonempty family $\mathcal X$ of sets, let
\[
\mathit{choose}(\mathcal X)\in\mathcal X
\]
denote an arbitrary admissible member. Thus, $\mathit{choose}$ represents
the nondeterministic choice made independently at each application of the
transition rule when constructing the BPMN state space; its result is itself
a set of flows.

The flow-passing function
$\mathit{Pass}: F_B^{exec} \rightarrow \mathcal{P}(F_B)$ is defined by:
\[
\mathit{Pass}(f) =
\left\{
\begin{array}{ll}
\{ f \}, & \text{if } t \in A \cup \StartE \cup \EndE \\[4pt]
\{ f \}, & \text{if } \gtype(t) = \mathit{Exclusive} \\[4pt]
\InF(t), & \text{if } \gtype(t) = \mathit{Parallel} \\[4pt]
\mathit{choose}\!\left(
\{\,S\subseteq\InF(t)\mid f\in S\,\}
\right), & \text{if } \gtype(t) = \mathit{Inclusive}
\end{array}
\right.
\]

The successor-selection function
$\mathit{Next}: F_B^{exec} \rightarrow \mathcal{P}(F)$ is defined by:
\[
\mathit{Next}(f) =
\left\{
\begin{array}{ll}
\OutF(t), & \text{if } t \in A \cup \StartE \\[4pt]
\varnothing, & \text{if } t \in \EndE \\[4pt]
\mathit{choose}\!\left(
\{\,S\subseteq\OutF(t)\mid |S|=1\,\}
\right), & \text{if } \gtype(t) = \mathit{Exclusive} \\[4pt]
\mathit{choose}\!\left(
\{\,S\subseteq\OutF(t)\mid S\neq\varnothing\,\}
\right), & \text{if } \gtype(t) = \mathit{Inclusive} \\[4pt]
\OutF(t), & \text{if } \gtype(t) = \mathit{Parallel}
\end{array}
\right.
\]

Consequently, executing a flow $f=(u,e)$ with $e\in\EndE$ consumes $f$
and produces no successor flow. The End Event has no precondition, but the
branch-specific postcondition contributed by its source $u$ is still checked
through $\mathit{Post}_B(f)$ in the validation semantics below. Since
$\varnothing\subseteq F_B$ and the final-state predicate is universal, the
empty flow configuration obtained after all terminal flows have been consumed
is also final.


\section{BPMN Validation}

OCL constrains this pair but neither constructs $\sigma_M'$ nor performs the state change.

Let the node precondition and the branch-specific postcondition be given by
\[
\mathit{pre}_T:
T\setminus\EndE\rightarrow\mathit{BoolExpr}(M),
\qquad
\mathit{post}_F:
F\rightarrow\mathit{BoolPostExpr}(M).
\]

For a flow $f=(u,v)$, its precondition is taken from its target $v$, whereas
its postcondition is taken from its source $u$ and is specific to the branch
$f$. Define the flow contracts
\[
\mathit{Pre}_B(f)=
\left\{
\begin{array}{ll}
\mathit{pre}_T(\mathit{target}(f)),
& \text{if }\mathit{target}(f)\notin\EndE,\\[2mm]
\top,
& \text{if }\mathit{target}(f)\in\EndE,
\end{array}
\right.
\]
and
\[
\mathit{Post}_B(f)=
\left\{
\begin{array}{ll}
\top,
& \text{if } f\in F^{0},\\[2mm]
\mathit{post}_F(f),
& \text{if } f\in F.
\end{array}
\right.
\]
Hence no condition is evaluated merely because a flow is selected by
$\mathit{Pass}$ or $\mathit{Next}$. The two contracts attached to $f$ are
evaluated only when $f$ occurs as an execution step.

Let $\mathcal{R}(B)$ denote the set of all finite executions reachable
from the initial configuration of $B$.

For

\[
\rho_B =
\sigma_F^0
\xrightarrow{\mathit{target}(f_1)}
\sigma_F^1
\cdots
\xrightarrow{\mathit{target}(f_n)}
\sigma_F^n,
\qquad
f_i \in \sigma_F^{i-1},
\qquad
1 \leq i \leq n,
\]

an ACL state path corresponding to $\rho_B$ is a sequence

\[
P_M =
\langle
\sigma_M^0,
\sigma_M^1,
\ldots,
\sigma_M^n
\rangle,
\qquad
\sigma_M^i \in \Sigma(M).
\]

We write

\[
P_M \models_M \rho_B
\]

iff

\[
\forall i\in\{0,\ldots,n\}:
\quad
\sigma_M^i \models Inv(M),
\]

\[
\forall i\in\{1,\ldots,n\}:
\quad
\sigma_M^{i-1} \models \mathit{Pre}_B(f_i),
\]

and

\[
\forall i\in\{1,\ldots,n\}:
\quad
(\sigma_M^{i-1},\sigma_M^i)
\models \mathit{Post}_B(f_i).
\]

The BPMN process $B$ is valid with respect to $M$ iff

\[
Valid_M(B)
\iff
\forall \rho_B\in\mathcal{R}(B)\;
\exists P_M\in\Sigma(M)^{n+1}:
P_M\models_M\rho_B.
\tag{VALID}
\]

\subsection{Bounded relational validation}

For an executable finite check, let the cyclic-flow set be
\[
\mathit{Cyc}(B)
=
\{\,(u,v)\in F\mid u\text{ is reachable from }v\text{ through }F\,\}.
\]
For a non-negative bound $k$, define
\[
\mathcal R_k(B)
=
\{\,\rho\in\mathcal R(B)
\mid
\forall f\in\mathit{Cyc}(B):\mathit{occ}_f(\rho)\leq k\,\}.
\]
Let a boundary $\beta$ assign a finite upper and lower object scope to every
ACL classifier, finite value domains to primitive types, a maximum number of
snapshots, and the cyclic-flow bound $k$. It induces
\[
D_\beta
=
\{\,\sigma_M\in\Sigma(M)\mid\sigma_M\text{ satisfies the scopes of }\beta\,\}.
\]
If $s_\beta$ is the maximum number of snapshots, define
\[
\mathcal R_{k,s_\beta}(B)
=
\{\,\rho\in\mathcal R_k(B)\mid |\rho|+1\leq s_\beta\,\}.
\]
Let
$\mathit{Max}(\mathcal R_{k,s_\beta}(B))$ contain the executions that are
maximal within both bounds (all flows have been consumed, or every remaining
extension would exceed $k$ or $s_\beta$). The bounded validation problem is
\[
Valid_M^{k,D_\beta}(B)
\iff
\forall\rho_B\in\mathit{Max}(\mathcal R_{k,s_\beta}(B))\;
\exists P_M\in D_\beta^{n+1}:
P_M\models_M\rho_B.
\tag{BVALID}
\]
Checking only maximal bounded executions is sufficient: a witness path for a
maximal execution restricts to a witness for each of its prefixes.

The executable validator enumerates the universal BPMN choices from
\textsc{Exec}. For each resulting $\rho_B$, Kodkod introduces a symbolic copy
of the ACL state relations at every time index:
\[
\sigma_M^i
=
(\sigma_{Class}^i,\sigma_{Att}^i,\sigma_{Assoc}^i,\sigma_{Play}^i).
\]
Their bounds are derived solely from $\beta$; no concrete object diagram is
supplied. The SAT query is
\[
\bigwedge_{i=0}^{n}
\bigl(WF_A(\sigma_M^i)\land Inv_M(\sigma_M^i)\bigr)
\quad\land\quad
\bigwedge_{i=1}^{n}
\left(
\mathit{Pre}_B(f_i)(\sigma_M^{i-1})
\land
\mathit{Post}_B(f_i)(\sigma_M^{i-1},\sigma_M^i)
\right).
\]
SAT supplies the existential path $P_M$; UNSAT supplies a bounded BPMN
execution for which no object-diagram path exists in $D_\beta$. Consequently,
$Valid_M^{k,D_\beta}(B)$ is a bounded proof over all states admitted by
$\beta$, not a proof of
$Valid_M(B)$ over unenumerated elements of $\Sigma(M)$.


\end{document}
